% theme: meiji_restoration_to_russo_japanese_war
\MajorQuestion{明治維新〜日清・日露戦争}
\SetRowSpaceQanda

\Instruction{明治新政府の成立に関する次の問いに答えなさい。}
\QQ{明治新政府の基本方針を示した文書を何というか。}{五箇条の御誓文}
\QQ{藩を廃止し、府と県を置いた改革を何というか。}{廃藩置県}
\QQ{江戸時代の身分制度を改めた政策を何というか。}{四民平等}
\QQ{欧米の制度や文化を調査するために派遣された使節団を何というか。}{岩倉使節団}
\QQ{北海道の開拓や警備のために置かれた役所を何というか。}{開拓使}

\Instruction{明治政府の改革に関する次の問いに答えなさい。}
\QQ{6歳以上の男女に小学校教育を受けさせる制度を何というか。}{学制}
\QQ{満20歳以上の男子に兵役を義務づけた法令を何というか。}{徴兵令}
\QQ{土地の所有者に地券を発行し、地価の3\%を現金で納めさせた改革を何というか。}{地租改正}
\QQ{経済力をつけ、強い軍隊をもつことを目指す政策を何というか。}{富国強兵}
\QQ{官営模範工場をつくり、近代産業を育成した政策を何というか。}{殖産興業}

\Instruction{自由民権運動と憲法に関する次の問いに答えなさい。}
\QQ{議会政治の実現を目指した運動を\NamedBlank{freedom_rights_movement}という。}{自由民権運動}
\QQ{板垣退助らが国会の開設を求めて提出した意見書を何というか。}{民撰議院設立の建白書}
\QQ{板垣退助を中心に結成された政党を何というか。}{自由党}
\QQ{大隈重信を中心に結成された政党を何というか。}{立憲改進党}
\QQ{伊藤博文らがドイツの憲法をもとに起草した憲法を何というか。}{大日本帝国憲法}

\Instruction{日清・日露戦争と近代日本に関する次の問いに答えなさい。}
\QQ{朝鮮をめぐる日本と清の対立から始まった戦争を何というか。}{日清戦争}
\QQ{\RefBlank{freedom_rights_movement}の後、日清戦争で日本と清の間に結ばれた条約を何というか。}{下関条約}
\QQ{ロシア・ドイツ・フランスが日本に遼東半島の返還を求めたことを何というか。}{三国干渉}
\QQ{日露戦争の講和条約を何というか。}{ポーツマス条約}
\QQ{日清戦争前後に軽工業、日露戦争前後に重工業が発達した。この変化を何というか。}{産業革命}
